\section{RAG System}
\label{sec:rag}

\subsection{Three-Mode Architecture}

The RAG system exposes three query modes of increasing capability, connected by a cascading fallback chain to ensure reliable answers under all conditions.

\begin{table}[H]
    \centering
    \caption{RAG Query Modes}
    \begin{tabular}{@{}llll@{}}
        \toprule
        \textbf{Mode} & \textbf{Strategy} & \textbf{Fallback Trigger} & \textbf{Fallback To} \\ \midrule
        LLM Direct   & Pure LLM, no retrieval          & ---                              & --- \\
        Fast RAG     & Single-pass vector retrieval    & Relevance $<$ 0.3 or $<$ 2 docs & LLM Direct \\
        Agentic RAG  & Multi-step ReAct agent          & Agent failure / low quality    & Fast RAG $\to$ LLM Direct \\ \bottomrule
    \end{tabular}
\end{table}

The fallback chain ensures the system always returns a useful response. Fast RAG falls back when vector retrieval cannot find sufficiently relevant documents (minimum relevance threshold \texttt{MIN\_RELEVANCE\_THRESHOLD=0.3}, minimum document count \texttt{MIN\_DOCUMENTS\_REQUIRED=2}). Agentic RAG falls back to Fast RAG on agent-level failures and ultimately to LLM Direct if retrieval also fails.

\subsection{Vector Database Design}

\subsubsection{ChromaDB Selection}

We selected ChromaDB as our vector database for several reasons:

\begin{itemize}
    \item \textbf{Embedded Operation:} No separate server required; runs within the Django process
    \item \textbf{Persistent Storage:} Data survives process restarts via SQLite backend
    \item \textbf{Metadata Filtering:} Supports filtering queries by metadata fields (video\_id, type, etc.)
    \item \textbf{HNSW Indexing:} Efficient approximate nearest neighbor search using Hierarchical Navigable Small Worlds algorithm
    \item \textbf{Simple API:} Clean Python interface for upsert and query operations
\end{itemize}

\subsubsection{Embedding Model}

We use \texttt{all-MiniLM-L6-v2} from sentence-transformers:

\begin{table}[H]
    \centering
    \caption{Embedding Model Specifications}
    \begin{tabular}{@{}ll@{}}
        \toprule
        \textbf{Property} & \textbf{Value} \\ \midrule
        Architecture & 6-layer MiniLM (distilled BERT) \\
        Output dimensions & 384 \\
        Model size & ~80 MB \\
        Max sequence length & 256 tokens \\
        Encoding speed (CPU) & ~100 texts/sec \\
        Memory footprint & ~200 MB \\ \bottomrule
    \end{tabular}
\end{table}

This model provides an excellent balance of performance and efficiency for our use case. While larger models (e.g., all-mpnet-base-v2) offer marginally better quality, the MiniLM model is sufficient for educational content and significantly faster.

\subsubsection{Document Types}

Three types of content are embedded and stored:

\textbf{1. Knowledge Points} (type: \texttt{knowledge\_point})

Embed text format:
\begin{verbatim}
"{title}: {summary} (Key terms: {term1}, {term2})"
\end{verbatim}

Metadata includes video\_id, section\_id, title, begin\_time, end\_time, and importance score.

\textbf{2. Section Transcripts} (type: \texttt{section\_transcript})

Embed text: First 2000 characters of section transcript (truncated to fit model limits).

Metadata includes video\_id, section\_id, title, begin\_time, end\_time.

\textbf{3. Slide OCR Text} (type: \texttt{slide\_ocr})

Embed text: OCR-extracted text content from high-resolution slide thumbnails (1920px), enabling semantic search over slide textual content including equations, bullet points, and diagram labels not captured in the spoken transcript.

Metadata includes video\_id, thumbnail\_id, timestamp.

\subsubsection{Storage Efficiency}

Typical storage per 60-minute lecture:

\begin{itemize}
    \item Knowledge points: 25-40 entries $\times$ ~200 bytes = 5-8 KB
    \item Section transcripts: 8-12 entries $\times$ ~2 KB = 16-24 KB
    \item Slide OCR entries: 20-60 entries $\times$ ~500 bytes = 10-30 KB
    \item HNSW index overhead: ~100-200 KB
    \item \textbf{Total:} 2-4 MB per video (including SQLite overhead)
\end{itemize}

\subsection{Fast RAG Pipeline}

The Fast RAG mode provides single-shot retrieval and generation with quality-gated fallback.

\subsubsection{Retrieval Process}

Given a user question $q$:

\begin{enumerate}
    \item Compute embedding: $e_q = E(q)$
    \item Query vector store with video\_id filter:
    \begin{equation}
        R = \text{top\_k}(\{d : \cos(e_q, e_d) \geq \tau_{rel}\})
    \end{equation}
    where $\tau_{rel} = 0.3$ (relevance threshold)
    \item If $|R| < 2$: fall back to LLM Direct mode
    \item Otherwise: format retrieved documents as numbered sources
\end{enumerate}

\subsubsection{Context Assembly}

Retrieved sources are formatted into a structured prompt:

\begin{lstlisting}[language=text, caption={RAG Context Template}]
## Lecture Context

### Video: {video_title}

### Lecture Overview
{knowledge_summary}

### Retrieved Sources:
[Source 1] (knowledge_point) "Learning Rate" [02:00-03:00] 
           (relevance: 0.85)
The learning rate controls the step size during gradient descent...

[Source 2] (section_transcript) "Optimization Techniques" [05:00-06:30]
           (relevance: 0.72)
There are several optimization techniques we can use...

---
Student Question: {question}
\end{lstlisting}

\subsubsection{Answer Generation}

The assembled prompt is sent to the LLM (qwen3-max) with the system prompt:

\begin{lstlisting}[language=text, caption={RAG System Prompt}]
You are a knowledgeable teaching assistant for a video lecture.
Answer the student's question based on the lecture content provided.

Instructions:
- Answer ONLY based on the provided context
- Cite sources using [Source N] notation
- Be concise but thorough, use markdown formatting
- Maintain an educational, helpful tone
\end{lstlisting}

\subsubsection{Citation Extraction}

After generation, citations are extracted by parsing [Source N] references and matching them to the source metadata:

\begin{lstlisting}[language=json, caption={Citation Structure}]
{
  "source_num": 1,
  "title": "Learning Rate",
  "begin_time": 120.0,
  "end_time": 180.0,
  "type": "knowledge_point",
  "relevance": 0.85
}
\end{lstlisting}

These citations are returned to the frontend for rendering as clickable timestamp badges.

\subsection{Agent System (ReAct Mode)}

For complex questions requiring multi-step reasoning, we implement a ReAct (Reasoning + Acting) agent using the LangGraph framework.

\subsubsection{ReAct Loop}

The agent follows an iterative loop implemented as a LangGraph state machine:

\begin{enumerate}
    \item LLM analyzes question and decides whether to call a tool or respond
    \item If tool call: execute tool, append result to conversation
    \item Repeat until LLM produces final text response
    \item Maximum 5 iterations to prevent infinite loops
    \item On failure or low quality: fall back to Fast RAG
\end{enumerate}

\subsubsection{Available Tools}

Six tools are available to the agent:

\begin{table}[H]
    \centering
    \caption{Agent Tools}
    \begin{tabular}{@{}lll@{}}
        \toprule
        \textbf{Tool} & \textbf{Purpose} & \textbf{Input} \\ \midrule
        search\_knowledge & Vector search on knowledge points \& transcripts & query string \\
        search\_slides    & OCR-text search on slide content                & query string \\
        get\_section\_details & Full section content & section order \\
        get\_lecture\_summary & Overview + chapters & --- \\
        list\_sections & All section titles/times & --- \\
        get\_transcript\_at\_time & Transcript slice & start, end times \\ \bottomrule
    \end{tabular}
\end{table}

The \texttt{search\_slides} tool searches \texttt{SlideOCR} records stored in ChromaDB, giving the agent access to text content written on slides (equations, definitions, bullet points) that may not be verbally explained in the transcript.

\subsubsection{Citation Sanitization}

A post-processing step strips hallucinated citation markers from agent responses. When the agent fabricates references such as \texttt{[Source 4]} that do not correspond to any actual tool result, these are removed from the final answer before delivery to the frontend.

\subsubsection{Tool Selection}

The LLM uses function calling to select tools. Each tool is described with a JSON schema:

\begin{lstlisting}[language=json, caption={Tool Schema Example}]
{
  "name": "search_knowledge",
  "description": "Search the lecture knowledge base for relevant content",
  "parameters": {
    "type": "object",
    "properties": {
      "query": {
        "type": "string",
        "description": "Search query describing the information needed"
      }
    },
    "required": ["query"]
  }
}
\end{lstlisting}

\subsubsection{Example Execution}

Question: "How does the lecture explain the relationship between learning rate and convergence?"

\begin{verbatim}
Iteration 1:
  LLM Decision: Call search_knowledge(query="learning rate convergence")
  Tool Result: [Result 1] (knowledge_point) "Learning Rate..." ...

Iteration 2:
  LLM Decision: Call get_section_details(section_order=3)
  Tool Result: "## Section 3: Optimization Techniques..."

Iteration 3:
  LLM Decision: Compose final answer
  Output: "The lecture explains that the learning rate directly 
           affects convergence speed..."
\end{verbatim}

\subsubsection{Streaming Protocol}

The agent uses Server-Sent Events (SSE) to stream progress:

\begin{itemize}
    \item \texttt{thinking}: Reasoning steps ("Analyzing question...")
    \item \texttt{tool\_call}: Tool invocations with arguments
    \item \texttt{tool\_result}: Tool outputs (300-char preview)
    \item \texttt{token}: Answer text chunks (streaming)
    \item \texttt{citations}: Structured citation list
    \item \texttt{done}: Completion signal with tool history
\end{itemize}

This provides immediate feedback to users during the 5-15 second processing time.

\subsection{Multi-Turn Conversations}

For follow-up questions, conversation history is included in the prompt:

\begin{itemize}
    \item Last 6 messages (3 turns) are included
    \item Provides context for pronouns ("explain that more")
    \item Enables clarifying questions and refinements
    \item Session persisted in database for later resumption
\end{itemize}

The 6-message limit balances context quality with token budget constraints.

